\section{Communion and Conversion: Developed Commentary}
\zlabel{triptych:concise:commentary:start}

\subsection{A complete people whose members still need mercy}
John's vision follows the question of who can stand before divine wrath
at the end of Apocalypse 6. In chapter 7 the servants of God are sealed
before the threatened injury proceeds. The appointed Lesson begins at
verse 2, with the angel who commands that delay. The multitude later
stands before the throne and the Lamb. Belonging to God precedes the
description of its worship. The definite number in between requires an
explanation: the twelve tribes are named, with twelve thousand repeated
for each, before a multitude appears which nobody can count.

Bede reasons from the meaning of Abraham's descent. Paul calls those who
belong to Christ Abraham's seed; the patriarchal inheritance therefore
admits imitation of faith as well as natural descent. Bede then explains
the number through the twelve apostles sent throughout the world.
Multiplication expresses the Church's completeness. On this account,
the exactness of 144,000 signifies a complete people, without supplying
a maximum number beyond which salvation becomes impossible. His account
uses the tribal list itself, rather than passing it over to reach the
multinational scene.\footnote{Apoc.~6:12--7:12; Gal.~3:29; Bede,
\work{Explanation}, I, on Apoc.~7:4, p.~46.}

The alternative Bede permits on verse 9 preserves a more immediate
distinction between the groups: Israel, to whom the Gospel was first
preached, is followed by the Gentiles' salvation. He also explains
the vision's return to victory after a recapitulation. These explanations
share the admission of the nations, but a reader cannot substitute
one for the other at every point. One takes the numbered company as
the whole Church; the other distinguishes Israel and Gentiles in the
sequence. The first interpretation of this Mass can retain both
without suppressing John's repeated names or restricting the feast
to one earthly people.\footnote{Bede, \work{Explanation}, I,
on Apoc.~7:9, p.~52.}

Bede also understands the ordered tribal names as virtues, beginning with
Juda as confession and proceeding toward inheritance. This is his spiritual
exposition of the names, rather than an independent account of their
Hebrew derivation. The people he describes are formed in holiness as well
as complete in number. His moral explanation can therefore be compared
with the Gospel's formation of the blessed, though the two biblical lists
remain distinct.\footnote{Bede, \work{Explanation}, I,
on Apoc.~7:5--8, pp.~46--51.}

Augustine answers a different question about the extent of this people.
In \work{The City of God} XX.9 he discusses souls reigning with Christ
in Apocalypse 20. The Church remembers the departed at the altar when
Christ's body is received. For Augustine that practice has a reason:
the faithful dead remain members of Christ. Bodily death has changed
their condition while leaving that membership intact. The coming
resurrection will restore the bodies which their souls await.
\footnote{Augustine, \work{City of God}, XX.9. His direct
Apocalypse discussion concerns chapter 20.}

Read beside the feast's Collect and Secret, this account explains
the living's relation to those they honor. The Collect asks mercy
amid many intercessors; the Secret asks a saving benefit for the
offerers. The saints' prayers are stated by the orations themselves.
Augustine explains continued membership, and that explanation helps
us understand how death leaves the saints within the communion
of those offering and praying. Neither Bede's symbolic total nor
Augustine's account makes the living assembly already identical in
holiness with the saints it venerates.

Augustine explicitly distinguishes the present Church, which contains
wheat and tares, from the kingdom finally free of evil. The Sunday
Collect's confession of sins therefore fits the first interpretation's
account of belonging. God's people ask release from what their frailty
has contracted. In the second interpretation, the same petition
concerns the sinful attachments which conversion must loosen.
The first explains why the people can ask with the saints; the second
explains the change for which the people ask.

\subsection{Holiness received and the labor of learning}
Bede's explanation of the multitude also concerns how it became holy.
He identifies its white robes with baptism and its palms with the
victory of the Cross. He permits a further account of the robes as
glory given by the Holy Spirit. At the final praise he says that the
Church acknowledges in each member what it has received from God.
The robes and acclamation describe people endowed with gifts,
and the palms ascribe their victory to Christ.\footnote{Bede,
\work{Explanation}, I, on Apoc.~7:9--12, pp.~52--53.}

In the first interpretation, this explains the direction of praise
in the entrance chant. The saints are honored while the angels
praise the Son. Schuster's account of sanctity in every condition
of life concerns members of one body sanctified by the Spirit.
Different offices, social positions, and languages remain compatible
with their common thanksgiving. Augustine's exposition of the
Introit's Psalm verse adds a condition on the singers: upright
praise submits to God. Rejoicing in him requires acceptance of
his goodness by the person who praises.\footnote{Schuster,
\work{Sacramentary}, V, pp.~208--210; Augustine,
Ps.~32, enarratio I, \S1.}

Read alongside Christ's invitation to the burdened, Bede's baptismal
robes suggest a further consequence. The person who labors at conversion
responds to a holiness already received. The Sunday Secret asks God
to complete an unmerited gift and increase service. Schuster explains
that service through ministry at the altar and the sacrifice of praise.
The work asked of the faithful remains dependent on what God gives.
The feast's Secret likewise requests the saving fruit of offerings
through mercy. These petitions place both beginning and perseverance
within dependence on God.\footnote{MR 1962, pp.~417, 703,
Secrets; Schuster, \work{Sacramentary}, III, p.~188.}

Augustine's account of the Beatitudes describes the work of that
conversion in detail. Humility makes a person teachable. Meekness
accepts Scripture even when an unfamiliar passage is not understood.
Mourning follows the recognition of attachments which have held the
soul away from the highest good. The former delights are being lost;
love of eternal goods is still being formed. Hunger for justice
describes the effort to desire the good strongly enough to abandon
the old delight. The sequence explains why sorrow can accompany
conversion even when its end is blessedness.\footnote{Augustine,
\work{Sermon on the Mount}, I.1.3; I.2.5--6; I.3.10.}

Mercy then concerns the person who needs help and can also give it.
Augustine reasons that someone seeking assistance from one stronger
should assist another who is weaker. The need to receive mercy
creates a reason to show it. Read beside the Collect's request,
this gives the petitioners a concrete disposition: those asking
mercy relieve another's misery. Their dependence does not suspend
their responsibility for other people.\footnote{Augustine,
\work{Sermon on the Mount}, I.2.7; I.3.10.}

The two interpretations consequently reach different conclusions
from the gifts celebrated at All Saints. Membership explains
common praise and mutual assistance across different conditions.
The formation of desire explains why those who have received a gift
still need instruction, relinquishment, and practical mercy.
The commemorated Sunday's request for increased service belongs
beside both: service is an act within Christ's people and a
continuing change in the person who offers it.

\subsection{Rest, hunger, and the goods Christ promises}
The Alleluia's invitation raises the second interpretation's hardest
question. Christ promises refreshment, yet the Gospel blesses mourning,
hunger for justice, and persecution. Chrysostom answers through
Matthew 11:29--30, the context following the appointed verse.
Disciples learn Christ's meekness and humility and take his yoke.
Chrysostom recalls that Christ also began the Beatitudes with
humility. The same disposition explains both the blessed life
on the mountain and the ease of the yoke.\footnote{Matt.~11:25--30;
Chrysostom, \work{Matthew}, XXXVIII.3.}

His reasoning compares obedience with the burden which sin already
imposes. Pride requires continual satisfaction. Anger prolongs
a contest by retaliation. Anxiety over possessions multiplies
what a person fears to lose. Meekness and self-command can relieve
those particular burdens: declining to retaliate can end the
quarrel which another blow would prolong. The relief concerns a
changed will and its actions. It leaves the real cost of restraint,
loss, and fidelity visible.\footnote{Chrysostom,
\work{Matthew}, XXXVIII.3--5.}

Chrysostom himself asks how this account agrees with the cross
and the narrow way. He invokes Christ's example, the endurance
of apostles and martyrs, and glory greater than present suffering.
He insists on both terms of Christ's saying: disciples hear of
a yoke so that they expect labor, and of its ease so that they
will undertake it. Present freedom from sin's bondage and the
future reward are distinct parts of his answer. Reducing the
promise to either painless living or postponed relief would
lose part of his explanation.

The Gradual poses a related difficulty about the good God gives.
Augustine imagines the worshipper who seeks God but remains poor
while a wealthy wrongdoer dies with an expensive funeral.
The facts require an account of what the Psalm promises.
Augustine names justice, truth, charity, faith, and endurance
as riches of the heart. He then identifies Christ as living
bread and explicitly cites the hunger and thirst for justice
of the Beatitudes. The good sought can therefore be present
in someone who lacks the rich person's possessions.\footnote{Augustine,
Ps.~33, enarratio II, \S\S14--15; NPNF, Psalm XXXIV,
\S\S13--14. The chant omits the first clause of Ps.~33:11,
whose rich and hungry remain part of the Psalm's context.}

Aquinas retains several explanations of the same promise.
Spiritual treasure, contentment with enough, providential help,
and a physician's withholding of desired food concern different
ways a good is given. The physician comparison explains a
deprivation through the patient's benefit. These alternatives
leave room for bodily provision and for prayer in actual need;
they furnish no assurance that every desired possession will
be obtained.\footnote{Aquinas, \work{Super Psalmos}, Ps.~33,
nn.~10--11.}

The inherited earth reveals a real difference between the Fathers.
Augustine identifies an enduring land of the living, where the
saints rest with an inheritance from which they cannot be driven.
Chrysostom defends a sensible present reward. Meekness can preserve
possessions which rashness and anger would destroy. He also
retains future gifts. Augustine's enduring inheritance and
Chrysostom's present earth are different referents, though both
connect meekness with a good received from God.\footnote{Augustine,
\work{Sermon on the Mount}, I.2.4; Chrysostom,
\work{Matthew}, XV.5.}

\subsection{The pure heart and the reconciliation of others}
The Communion repeats three Gospel promises to those approaching
the sacrament: the pure see God, peacemakers are called God's
children, and those persecuted for justice possess the kingdom.
For the first interpretation, these are promises to members
of the same people whose departed faithful are honored.
For the second, their order describes the holiness still being
formed in those who hear them. Purity and peace must therefore
be understood closely enough to explain how they agree with
the persecution named immediately afterward.

Augustine compares the pure heart to an eye fitted for its object.
Bodily sight requires sound eyes; the sight of God requires
purification of the inward faculty suited to him. The pure
heart is single. This explanation follows the earlier work
of giving up contrary attachments. A person's desire has
to be directed toward the God whom that person hopes to see.
\footnote{Augustine, \work{Sermon on the Mount}, I.2.8.}

Chrysostom specifies a practical difficulty which a general
praise of generosity could leave unaddressed. Someone can
give alms freely, avoid greed, and still commit sins against
chastity. The blessing of purity requires more than that
one conspicuous virtue. Beside Augustine's singleness, his
example explains why mercy cannot excuse the retention of
an incompatible habit. His account of mercy already includes
actions beyond giving money.\footnote{Chrysostom,
\work{Matthew}, XV.6.}

Their explanations of peace are complementary. Augustine
describes order within the person: desires obey reason, and
reason itself obeys the truth, the only-begotten Son. The
highest human faculty needs submission to one higher.
This gives a reason why external hostility need not destroy
peace: the hostile act need not make the disciple's own
desires rebel against Christ. Augustine understands the
following persecution as testing the strength of what has
been established within.\footnote{Augustine,
\work{Sermon on the Mount}, I.2.9--I.3.10.}

Chrysostom requires the reconciliation of others who are
divided. He explains the title of God's children through
the work of the only-begotten Son, who reconciles the
alienated. Peace thus has a task in other people's relations,
as well as within one's own soul. He then qualifies it
through the next blessing: agreement in wrongdoing is
excluded. The persecution must be for justice; the subsequent
reviling must be false and for Christ. Accusation alone
does not establish a person's holiness.\footnote{Chrysostom,
\work{Matthew}, XV.7.}

Inward order and outward reconciliation explain different
aspects of the same promised sonship. They also explain
why the feast can repeat peace and persecution together.
Reconciliation remains a duty when it is possible in
justice. Fidelity remains a duty when an offered agreement
requires evil. The final petitions for protection belong
to disciples still learning that fidelity after sharing
in divine gifts.

\subsection{Peace after death and fulfillment still awaited}
The Offertory's Wisdom selection concerns the judgment of
those who think the righteous have perished. Wisdom 2
describes persecutors testing the just man's claim to
be God's child. Chapter 3 answers their assessment of
his death. The chant selects and adapts its clauses,
using the torment of malice where the ordinary biblical
wording speaks of death. It proclaims the just souls'
security and peace without denying their bodily killing.
\footnote{Wis.~2:12--3:9; MR 1962, p.~703, Offertory.}

Augustine invokes Wisdom 3:1 in his exposition of Psalm 87
while explaining the righteous dead's condition. Their
bodies await resurrection; their souls remain in God's
hand. Christ freely lays down his life and takes it up;
the members receive bodily restoration after their Head.
The continued protection explains the peace of those
still awaiting that restoration. Together with
\work{The City of God} XX.9, this gives the first
interpretation both continued membership and a precise
account of what death has left incomplete.\footnote{Augustine,
Ps.~87, \S5 (NPNF, Psalm LXXXVIII, \S4);
\work{City of God}, XX.9.}

Gregory develops another consequence of suffering in his
explanation of Job's tree which grows again after being
cut down. He directly quotes Wisdom 3:2 when describing
despised preaching and a righteous person's killed body.
The Spirit makes that example awaken virtue in others.
In the following paragraph, Gregory gives a distinct
account of renewed growth as bodily resurrection.
One explanation concerns fruit in the living; the other
concerns restoration of the sufferer's body. His Wisdom
reuse consequently illuminates both conversion and
resurrection without identifying those two results.
\footnote{Gregory, \work{Moralia}, XII.5.6--7,
on Job~14:8--9, directly invoking Wis.~3:2.}

The difference matters to the four senses. The literal
account acknowledges persecution and bodily death.
Fulfillment in Christ concerns his victory and the members
who belong to him. The moral consequence includes the
virtue awakened by faithful endurance. Final fulfillment
includes the body's resurrection, which an account of
present peace cannot replace. Augustine also allows
present fulfillment of the Beatitudes while preserving
a future transformation beyond expression.\footnote{Augustine,
\work{Sermon on the Mount}, I.4.12.}

The heavenly praise likewise has distinct explanations.
Bede permits the angels to represent the worshipping
multitude, or to be actual heavenly spirits rejoicing
at Gentile salvation. The Introit names actual angels,
and the Common Preface names heavenly powers before
asking admission for human voices. Bede's latter
explanation can be compared directly with these prayers;
his symbolic explanation of the vision remains an
alternative at its own locus.\footnote{Bede,
\work{Explanation}, I, on Apoc.~7:11, p.~53;
MR 1962, pp.~252--253, 702.}

The first interpretation concludes with a people united
in Christ whose members' conditions differ: the living
ask mercy with the saints' help, and the departed
remain secure while awaiting resurrection. The second
concludes with the conversion Christ's rest permits
within suffering: humble desire, purity, and practical
peace anticipate the sight of God. In the feast's
Postcommunion the faithful seek enduring joy and the
saints' continual supplication. The Sunday petition
asks protection for those admitted to divine
participation. Schuster explains it as strengthening
within battle. The communion celebrated and the
conversion still sought meet in people asking for
the grace to persevere.\footnote{MR 1962, pp.~417, 703,
Postcommunions; Schuster, \work{Sacramentary}, III,
p.~189; V, p.~211.}
